% Analysis notes: crash_horizon_test_2
% Run: THE_CRASH, 5 LLM firms, 50 CES consumers, 100 max timesteps, unit cost from pref gen.
% Use state_t*.json and experiment_args.json from logs/crash_horizon_test_2/ for figures/numbers.

\documentclass{article}
\usepackage{microtype}
\usepackage{hyperref}
\usepackage{amsmath}
\usepackage{amssymb}
\usepackage{booktabs}

\title{Notes: Analysis of \texttt{crash\_horizon\_test\_2}}
\author{Agent Bazaar}
\date{}

\begin{document}
\maketitle

\section{Why Did Firms Start with Largely Negative Profits?}

In \texttt{crash\_horizon\_test\_2}, firms begin with \emph{step-level} profits that are largely negative and then creep toward zero. Two mechanisms explain this.

\textbf{(1) Fixed costs before revenue.} Each timestep, firms pay \textbf{overhead} (scaled to the daily timestep, e.g.\ \(\approx 50/7\) per day) and may pay \textbf{supply costs} when they purchase inputs. In early steps, firms set prices and production; if they sell little or nothing in that same step, revenue is small while overhead (and possibly supply expense) is still incurred. Thus the \emph{per-step} profit
\[
r_t^i = (P_t^i - c)\, Q_t^{i,\text{sold}} - \text{(overhead + supply + taxes + fees)}
\]
can be negative simply because fixed overhead (and any supply purchases) exceed revenue in that period. So ``profits started negative'' is consistent with the environment: fixed costs are charged every step regardless of sales.

\textbf{(2) Profits creeping to zero.} Step-level profit is reported only for \emph{active} firms. Once a firm goes \emph{bankrupt} (\(C_t^i < 0\)), it exits: it no longer sets prices, produces, or pays overhead, and its reported profit is typically zero (or the last step’s loss before exit). So as firms fail one by one, the cross-section of \emph{reported} profits is increasingly made up of zeros (failed firms) and possibly small or negative values for the last survivors. That gives the impression that ``profits crept to near 0''---in part because the number of firms still generating nonzero profit shrinks, and in part because the remaining firms may be pricing at or below unit cost and barely covering or not covering costs before they too fail.

\textbf{Summary.} Negative profits at the start are driven by fixed overhead (and supply costs) relative to low or zero revenue in early steps. The drift toward zero reflects both the exit of firms (so their profit is no longer reported) and the unsustainable pricing of survivors, consistent with the Crash dynamic.

\section{Firm Reputations and Their Effect on the Sim Outcome}

In the \textbf{B2C Crash} environment (THE\_CRASH), \textbf{reputation} is fulfillment-based: it is updated from the history of fulfilled vs.\ requested quantities (e.g., a rolling average of delivery performance), not from product quality as in the Lemon Market. In the Crash setup, \textbf{market clearing is by price}: consumers buy from the lowest-priced firm first (lexicographically by price, then firm index). So reputation does \emph{not} determine who gets demand; only prices and availability do.

Implications for \texttt{crash\_horizon\_test\_2}:
\begin{itemize}
    \item Reputations can remain \emph{high} (e.g., near 1.0) even as the market collapses, because firms may have fulfilled orders well until they ran out of inventory or went bankrupt.
    \item Reputation did \emph{not} prevent the crash: the failure mode is driven by prices falling below unit cost and cascading bankruptcies, not by a loss of buyer trust.
    \item If the sim records reputation over time, one may see reputations stay high until firms exit; the collapse is due to the \emph{price spiral}, not reputation decay. In the Lemon Market, by contrast, reputation is central and decay leads to volume collapse.
\end{itemize}

So in this run, firm reputations largely reflect that firms delivered what they sold; the outcome is dominated by the price war and solvency, not by reputation effects on demand.

\section{Price War and Its Relation to ``The Crash''}

The price war in \texttt{crash\_horizon\_test\_2} is the core of \textbf{Algorithmic Instability} (``The Crash'') as defined in the main paper (\texttt{root.pdf}).

\textbf{Mechanism.} Under \textbf{Search Friction}, each firm observes only a limited subset of the market (token budget \(\tau\)). Each firm acts on its local view: to attract demand, it has an incentive to undercut competitors. Because no firm can commit to a price floor and all are myopic, this leads to \emph{recursive undercutting}: \(\partial P_t / \partial t < 0\) until \(P_t^i < c\) for surviving firms. Once prices are below unit cost \(c\), every sale is loss-making; cash falls, and firms go bankrupt in sequence. The market \emph{collapses} when all firms are bankrupt (or effectively inactive).

\textbf{Why the 1-day timescale made the price war more aggressive.} With one timestep = one day, each period's demand and revenue are small (daily scale). Firms compete for that small slice every step and are incentivized to undercut ``just this period'' to capture today's sales. With many such steps, undercutting is repeated every day, so the spiral has more rounds and prices are driven below cost faster than under a coarser (e.g.\ weekly) timescale where each round has larger stakes and fewer rounds in the same horizon.

\textbf{Contextualization in the paper.} The abstract and Section~3.2 (Environment A: The Crash) state that in B2C markets, myopic firms amplify price volatility until the market collapses---an LLM-native analog of the 2010 Flash Crash. The reward is \(r_t^i = (P_t^i - c) Q_t^{i,\text{sold}} - f Q_t^i\); when \(P_t^i < c\), the margin term is negative and cash drains. Our run exhibits exactly this: a \textbf{price spiral} (price war), prices at or below unit cost, and \textbf{cascading bankruptcies}. Recovery metrics (e.g., timesteps to recovery after a supply shock) in the paper apply to experiments where a shock is introduced; in a baseline run like \texttt{crash\_horizon\_test\_2}, the instability is endogenous and the ``recovery'' is failure unless agents (e.g., Stabilizing Firms) maintain a price floor above \(c\).

\textbf{Report.} In \texttt{crash\_horizon\_test\_2}, the price war is the defining feature: average market price \(\bar{P}_t\) trends down, crosses below unit cost, and firms exit. This is the Crash failure mode; the run is a demonstration of why base (unaligned) LLM firms are economically unsafe and why finetuning for stability (Stabilizing Firms) is needed to keep \(\bar{P}_t \geq c\) and avoid collapse.

\section{Consumer Utility: Why Negative, and Is It Concerning?}

\subsection{Why Was Consumer Utility Negative?}

Consumer utility in the sim is
\[
U = U_{\text{goods}} + U_{\text{cash}} - D_{\text{labor}},
\]
where \(U_{\text{goods}}\) is CES utility from consumed goods, \(U_{\text{cash}} = \beta \ln(\varepsilon + \text{cash})\), and \(D_{\text{labor}} = c \cdot l^\delta\) (scaled for the timestep). So total utility can be negative when \textbf{labor disutility} exceeds the sum of goods and cash utility.

In \texttt{crash\_horizon\_test\_2}:
\begin{itemize}
    \item After the \textbf{market collapse}, many consumers receive \emph{zero} goods (no food purchased), so \(U_{\text{goods}} = 0\).
    \item \(U_{\text{cash}}\) is bounded (log of cash); it helps but does not dominate.
    \item \textbf{Labor} \(l\) is chosen by the consumer (e.g., 40--65 hours); \(D_{\text{labor}} = c \cdot l^\delta\) is convex in \(l\). For high \(l\) and \(\delta > 1\), \(D_{\text{labor}}\) becomes large.
\end{itemize}

So \textbf{negative utility} is expected when: (i) consumers work many hours (high \(l\)), (ii) they get little or no consumption (goods utility near 0 after the crash), and (iii) the convex labor disutility term dominates. The negative values in the stats are therefore consistent with the design: they indicate that, for those consumers, the disutility of labor was not offset by consumption and cash utility---especially once the market had collapsed and consumption dropped to zero.

\subsection{Cardinal vs Ordinal: Should We Be Concerned?}

The consumer utility function is \textbf{cardinal}: it uses a specific CES form, log cash, and a power-law labor disutility. Cardinal utility assigns meaning to the \emph{level} and \emph{magnitude} of utility (e.g., ``utility is \(-50\)''). \textbf{Ordinal} utility would only care about \emph{rankings} (prefer A to B); adding constants or applying monotone transforms would preserve ordinal content but change levels.

\textbf{Implications:}
\begin{itemize}
    \item \textbf{Ordinal content:} Comparisons across consumers or across time (e.g., ``consumer A is better off than B,'' or ``welfare improved after a policy change'') are valid in an ordinal sense if we only use the \emph{sign of differences}. Negative utility \emph{levels} do not by themselves imply an ordinal problem.
    \item \textbf{Cardinal content:} Our spec is cardinal, so the \emph{numbers} (e.g., \(-140\) vs \(-20\)) are sensitive to the chosen scale and functional form. Saying ``utility dropped by 100 units'' has meaning only relative to this cardinalization.
    \item \textbf{Should we be concerned?} About \textbf{negative levels} per se: not for ordinal welfare. What matters for welfare analysis is (i) \emph{relative} utility across options or time (ordinal), and (ii) that the \emph{mechanism} is correct (market collapse \(\Rightarrow\) no goods \(\Rightarrow\) low or negative utility for high-labor consumers). So the negative utilities are \textbf{not} a sign of a bug; they are a sign that the Crash harmed consumers (less consumption, same or higher labor disutility). For the paper’s \textbf{Consumer Welfare} dimension (\(\Phi_W\), aggregate consumer surplus or utility), we can use the same cardinal measure consistently before and after interventions; the \emph{comparison} (e.g., aligned agents improve \(\Phi_W\)) is then interpretable. We should \emph{not} over-interpret the absolute value of utility (e.g., ``-100 is twice as bad as -50'') unless we explicitly defend the cardinal scale.
\end{itemize}

\textbf{Conclusion.} Consumer utility was negative because labor disutility exceeded goods plus cash utility, especially after the market collapsed and consumption went to zero. The utility function is cardinal. Negative utility is not inherently concerning for ordinal welfare; it correctly reflects that welfare was poor when consumption collapsed. For the paper, report welfare in a consistent cardinal scale and emphasize comparisons (e.g., base vs.\ Stabilizing Firm runs) rather than the raw level of utility.

\textbf{Takeaways.} A smart firm could potentially hold out for a better price, halting production and sales until the price is at or above unit cost. This would prevent the price spiral and the market collapse. This would require firms to wait until average eWTP rebounds to above unit cost, which requires all firms to be aligned and to hold out for an extended period of time.

\end{document}
